\documentclass[10pt,a4paper]{amsart}
\usepackage[margin=19mm]{geometry}
\usepackage{amsmath,amssymb,mathtools}
\usepackage[scr=rsfs,cal=euler]{mathalfa}
\usepackage{xfrac}
\usepackage[unicode,hidelinks,bookmarks=false]{hyperref}
\newcommand{\ie}{i.e.\ }
\newcommand{\cf}{cf.\ }
\newcommand{\resp}{respectively\ }
\newcommand{\Subsection}{\subsection}
\providecommand{\nobreakdash}{\nobreak\hbox{-}}
\setlength{\emergencystretch}{2em}
\multlinegap0pt
\usepackage{amssymb}
\usepackage{xfrac}
\usepackage{mathtools}
\usepackage[shortlabels]{enumitem}
\usepackage{tikz}
\usepackage{xcolor}
\newtheorem{prop}{Proposition}
\newtheorem{theorem}{Theorem}
\usepackage{cleveref}
\crefname{prop}{Proposition}{Propositions}
\crefname{theorem}{Theorem}{Theorems}
\newcommand{\R}{\mathbb{R}}
\DeclareMathOperator{\besselspace}{H}
\newcommand{\dd}{\mathrm{d}} % This is the derivative d
\newcommand{\dispair}[2]{\langle #2,#1 \rangle}
\renewcommand{\vec}{\mathbf}
\title{The microscopic weighting on a metric space}
\author{Emily Roff, Simon Willerton}
\date{Source: arXiv:2607.05349v1 (CC BY 4.0)}
\begin{document}
\maketitle

\noindent\emph{Source and licence.} The microscopic weighting on a metric space. Version arXiv:2607.05349v1 (CC BY 4.0), distributed by arXiv under the Creative Commons Attribution 4.0 International licence, http://creativecommons.org/licenses/by/4.0/, as recorded in the arXiv licence metadata for this exact version and shown on the abstract page https://arxiv.org/abs/2607.05349v1. Full licence notice: \texttt{licenses/arxiv-2607.05349-CC-BY-4.0.txt}.\par\smallskip

\noindent\textit{Editorial scope.} Counted unit: the article's theorem thm:alexander\_limit (source0.tex lines 1592--1601). The article's complete original proof of it is delivered with it (source0.tex lines 1603--1702, one \texttt{proof} environment of 100 lines). No mathematics is written, repaired or re-derived: the statement and the proof are the article's own lines, in the article's own notation, and no retained line is substituted or re-flowed.  Retained uncounted as the source-local context the proof needs: lines 1484--1490; lines 1530--1545.  Nothing was omitted from any retained span, so the package carries no omission record.  No retained line contains a citation, so the wrapper carries no bibliography and no bibliography entry.  No retained line contains a figure environment, an image-inclusion command or drawing code, so no image is delivered and the package declares no asset dependency.  Editorial formatting only, in addition: an AMS \texttt{amsart} preamble that restates the article's own declarations and macros used by the retained lines, namely the environments \texttt{prop}, \texttt{theorem} on separate counters, together with the standard packages those lines need.  The retained environments print in this document as Proposition 1, Theorem 1 in order of appearance, whereas the article prints the same items with the numbering of its own full document.  The complete native source of the article as distributed in the arXiv e-print for this exact version is bundled unchanged as \texttt{sources/204536/source0.tex}.
\begin{equation}
\label{eq:Alexander-family}
    \nu_q \coloneq q \bar{s}_1 + (1 - q) \bar{s}_{1-1/q}
    \quad
    \text{for }
    q = 1, 2, \dots.
\end{equation}

\begin{prop}
\label{prop:microwtg-on-B3}
On \(B^3\), the microscopic weighting  \(\widehat{w} \in \besselspace^{-2}\) is given by
\[
    \dispair{\widehat{w}}{f}
    =
    \lim_{t \to 0}
    \dispair{w_t}{f}
    =
    \int_{S^2} f(\vec{s}) \,\dd \bar{s}
    +
    \int_{S^2} \tfrac{\partial}{\partial \vec{n}} f(\vec{s}) \,\dd \bar{s},
\]
for each function \(f \in \besselspace^2\). 
\qed
\end{prop}

\begin{theorem}
\label{thm:alexander_limit}
    Let \(\nu_1, \nu_2, \dots\) be Alexander's energy-maximizing family of measures on \(B^3\), and let  \(\widehat{w}\) be the microscopic weighting of \Cref{prop:microwtg-on-B3}. Let \(f \colon \R^3 \to \R\) be such that \(f\) and \(\tfrac{\partial}{\partial \vec{n}} f\) are both continuous on a neighbourhood of \(B^3\). Then
    \[
        \lim_{q \to \infty} \dispair{\nu_q}{f} = \lim_{q\to \infty} 
        \int_{B^3} f(\vec{b})  \,\dd \nu_q 
        =
        \dispair{\widehat{w}}{f}.
    \]
\end{theorem}

\begin{proof}
Using the definition of \(\nu_q\) in equation \eqref{eq:Alexander-family} and the expression for \(\widehat{w}\) in \Cref{prop:microwtg-on-B3}, we have
    \begin{align}
        \lim_{q\to \infty}
            \int_{B^3} f(\vec{b}) \,\dd \nu_q 
        \nonumber 
        &=
        \lim_{q\to \infty}
        \left(
            \int_{B^3} q f(\vec{b})  \,\dd \bar{s}_1 
            + 
            \int_{B^3} (1 - q)f(\vec{b}) \,\dd\bar{s}_{1-1/q}
        \right)
        \nonumber \\
        &=
        \lim_{q\to \infty}
        \left(
            \int_{S^2} q f(\vec{s})\,\dd \bar{s} 
            + 
            \int_{S^2} (1 - q) f\left(\left(1-\tfrac{1}{q}\right)\vec{s}\right)\,\dd\bar{s}
        \right)
        \label{eq:alexander_limit_1} \\
        &=
        \lim_{q\to \infty}
        \left(
            \int_{S^2} f\left(\left(1-\tfrac{1}{q}\right)\vec{s}\right) \,\dd \bar{s} 
            + 
            \int_{S^2} \frac{f(\vec{s}) - f\left(\left(1-\tfrac{1}{q}\right)\vec{s}\right)}{\tfrac{1}{q}} \,\dd\bar{s}
        \right)
        \label{eq:alexander_limit_2} \\
        &=
        \int_{S^2} \lim_{q\to \infty} f\left(\left(1-\tfrac{1}{q}\right)\vec{s}\right) \,\dd \bar{s} 
        + 
        \int_{S^2} \lim_{q\to \infty} \frac{f(\vec{s}) - f\left(\left(1-\tfrac{1}{q}\right)\vec{s}\right)}{\tfrac{1}{q}} \,\dd\bar{s}
        \label{eq:alexander_limit_3} \\
        &=
        \int_{S^2} f(\vec{s}) \,\dd \bar{s}
        +
        \int_{S^2} \tfrac{\partial}{\partial \vec{n}} f(\vec{s}) \,\dd \bar{s}
        \label{eq:alexander_limit_4} \\
        &=
        \dispair{\widehat{w}}{f}.
        \nonumber 
    \end{align}
Here, the equality \eqref{eq:alexander_limit_1} follows from a change of variables, using the fact that the measure \(\bar s_{1-1/q}\) is supported on an interior concentric sphere. The equality \eqref{eq:alexander_limit_2} is obtained by algebraic manipulation.

The exchange of integrals and limits in \eqref{eq:alexander_limit_3} is possible thanks to the dominated convergence theorem. As \(|f|\) is continuous on \(B^3\) it is bounded there, and if we choose \(g_1\) to be a smooth, compactly supported function that is constantly equal to \(\sup_{B^3}|f|\) on \(S^2\), then \(g_1\) is integrable and satisfies 
\[\bigl|f((1-1/q) \vec{s})\bigr| \leq g_1(\vec{s})\]
for all \(\vec{s} \in S^2\) and \(q 
\geq 1\). Similarly, we can choose \(g_2\) to be a smooth, compactly supported function that is constantly equal to \(\sup_{S^2} \left| \tfrac{\partial}{\partial \vec{n}} f \right| + 1\) on \(S^2\). We have 
\[
    \biggl|  
        \tfrac{1}{q} 
        \Bigl(
            f(\vec{s}) - f\bigl(\left(1 - 1/q \right)\vec{s}\bigr) 
        \Bigr)
    \biggr| 
    - 
    \biggl| 
        \tfrac{\partial}{\partial \vec{n}} f(\vec{s}) 
    \biggr| 
    \leq 
    \Biggl| 
        \biggl|  
            \tfrac{1}{q} 
            \Bigl(
                f(\vec{s}) - 
                f\bigl(
                    (1-1/q)\vec{s}
                \bigr) 
            \Bigr)
        \biggr| 
        - 
        \biggl| 
            \tfrac{\partial}{\partial \vec{n}} f(\vec{s}) 
        \biggr| 
    \Biggr| 
\]
and for \(q \gg 1\) the right hand side is bounded by \(1\), say, by the definition of the normal derivative.  Hence, for all \(\vec{s} \in S^2\) and \(q \gg 1\),
\[
    \biggl|  
        \tfrac{1}{q} 
        \Bigl(
            f(\vec{s}) - f\bigl(\left(1 - 1/q \right)\vec{s}\bigr) 
        \Bigr)
    \biggr| 
    < 
    1
    +
    \biggl| 
        \tfrac{\partial}{\partial \vec{n}} f(\vec{s}) 
    \biggr| 
    \leq 
    g_2(\vec{s})
    .
\]
Equality \eqref{eq:alexander_limit_3} follows from the dominated convergence theorem using \(g_1 + g_2\) as the dominating function.

Finally, the equality \eqref{eq:alexander_limit_4} follows using the continuity of \(f\) and the definition of the normal derivative.
\end{proof}

\end{document}
